\documentclass[11pt,a4paper]{article}
\usepackage{fontspec}
\usepackage[french]{babel}
\frenchsetup{AutoSpacePunctuation=false}
\usepackage[margin=2.0cm]{geometry}
\usepackage{amsmath,amssymb,amsthm}
\usepackage{enumitem}
\usepackage[colorlinks=true,linkcolor=blue!50!black,urlcolor=blue!50!black,citecolor=blue!50!black]{hyperref}
\usepackage{xcolor}
\setlength{\parskip}{3pt}
\setlength{\emergencystretch}{2em}

\newtheorem{theorem}{Théorème}
\newtheorem{proposition}[theorem]{Proposition}
\newtheorem{lemma}[theorem]{Lemme}
\newtheorem{corollary}[theorem]{Corollaire}
\theoremstyle{remark}
\newtheorem{remark}[theorem]{Remarque}

\newcommand{\Z}{\mathbb{Z}}
\newcommand{\Q}{\mathbb{Q}}
\newcommand{\F}{\mathbb{F}}
\newcommand{\one}{\mathbf{j}}
\newcommand{\srg}{\operatorname{srg}}
\newcommand{\Fix}{\operatorname{Fix}}
\newcommand{\tr}{\operatorname{tr}}
\newcommand{\rk}{\operatorname{rang}}

\title{Le graphe de Conway $\srg(99,14,1,2)$ n'a pas d'automorphisme d'ordre~3\\[3pt]
\large Une obstruction par les discriminants}
\author{Rédaction : Claude (Opus 5.5), pour D.~Djematene}
\date{23 septembre 2026 \\[4pt]
\textcolor{red!60!black}{\textbf{Preuve candidate, non encore relue par un tiers.}}\\[2pt]
\small Le théorème matriciel (théorème~5) est vérifié en Lean~4 / Mathlib (voir §5).}

\begin{document}
\maketitle

\begin{abstract}
Soit $\Gamma$ un graphe fortement régulier de paramètres $(99,14,1,2)$ et $g$ un automorphisme d'ordre~3 sans point fixe.
L'action de $g$ fournit deux matrices $33\times 33$ : le quotient $Q$, à coefficients entiers, et une matrice hermitienne $M$ à
coefficients dans les entiers d'Eisenstein $\Z[\omega]$, congrues modulo $1-\omega$. On montre qu'un tel couple ne peut pas exister :
les réseaux des vecteurs propres de valeur propre~$3$ de $Q$ et de $M$ ont des réductions modulo~3 égales, donc des
déterminants congrus modulo~3. Or le premier vaut $2$ modulo~3, à cause du facteur $11$ de $\langle\one,\one\rangle=33$,
et le second vaut $1$, car seul le premier $7$ y intervient. Le cas d'un triangle fixe se traite de la même façon.
Avec un théorème de Makhnev et Minakova, il en résulte que $\Gamma$ n'a aucun automorphisme d'ordre~3.
La preuve est de l'algèbre linéaire élémentaire sur $\Z$ et $\Z[\omega]$, sans calcul par ordinateur.
\end{abstract}

%====================================================================
\section{Énoncés}

Soit $A$ la matrice d'adjacence de $\Gamma=\srg(99,14,1,2)$ ; $A^2=14I+A+2(J-I-A)$, soit
\begin{equation}\label{eq:srg}
A^2+A=12I+2J ,
\end{equation}
où $J$ est la matrice dont tous les coefficients valent~$1$. On note $\one$ le vecteur dont toutes les coordonnées valent~$1$,
$\omega=e^{2i\pi/3}$ et $\pi=1-\omega$.

\begin{theorem}\label{thm:libre}
$\Gamma$ n'a aucun automorphisme d'ordre~$3$ sans point fixe.
\end{theorem}

\begin{proposition}\label{prop:triangle}
$\Gamma$ n'a aucun automorphisme d'ordre~$3$ dont l'ensemble des points fixes est un triangle.
\end{proposition}

\begin{corollary}\label{cor:tout}
$\Gamma$ n'a aucun automorphisme d'ordre~$3$.
\end{corollary}
\begin{proof}
D'après Makhnev et Minakova \cite{MM04}, les points fixes d'un automorphisme d'ordre~3 forment l'ensemble vide ou un triangle.
On conclut par le théorème~\ref{thm:libre} et la proposition~\ref{prop:triangle}. (Le cas avec points fixes est aussi exclu, par
une recherche sur ordinateur, dans \cite{B09,BL11}.)
\end{proof}

\emph{Contexte.} Makhnev et Minakova montrent que $|\mathrm{Aut}(\Gamma)|$ divise $2\cdot 3^3\cdot 7\cdot 11$ ; Cesarz et Woldar
\cite{CW23} montrent que si $2$ divise $|\mathrm{Aut}(\Gamma)|$, alors $|\mathrm{Aut}(\Gamma)|$ divise $6$. Le cas de l'ordre~3 sans
point fixe restait ouvert. L'identité trace--rang d'Ishida \cite{I26} y impose exactement six orbites triangulaires. L'argument
ci-dessous exploite l'invariant suivant, le \emph{discriminant}, et retrouve ce $\tau=6$ au passage (remarque~\ref{rem:tau}).

%====================================================================
\section{Réduction à un énoncé sur des matrices}

Soit $g$ un automorphisme d'ordre~3 sans point fixe. Ses $33$ orbites ont chacune trois éléments ; on choisit un représentant dans
chacune et l'on note les sommets $(U,b)$, $U\in\{1,\dots,33\}$, $b\in\Z/3$, avec $g(U,b)=(U,b+1)$. Pour $b\in\Z/3$, soit
$A_b\in M_{33}(\Z)$ définie par $A_b[U,V]=A[(U,0),(V,b)]$. Comme $g$ est un automorphisme, $A[(U,c),(V,c+b)]=A_b[U,V]$ pour tout
$c$, et la symétrie de $A$ donne $A_b^{\mathsf T}=A_{-b}$. On pose
\[
Q=A_0+A_1+A_2\in M_{33}(\Z),\qquad M=A_0+\omega A_1+\omega^2A_2\in M_{33}(\Z[\omega]).
\]

\begin{lemma}\label{lem:QM}
(i) $Q$ est symétrique, $Q\one=14\,\one$ et $Q^2+Q=12I+6J$.
(ii) $M$ est hermitienne et $M^2+M=12I$.
(iii) $Q-M=(1-\omega)A_1+(1-\omega^2)A_2$ est divisible par $\pi$ coefficient par coefficient.
(iv) $\tr Q=2\tau$ et $\tr M=-\tau$, où $\tau$ est le nombre d'orbites qui induisent un triangle.
\end{lemma}
\begin{proof}
Pour une matrice $X$ indexée par les sommets et commutant avec $g$, posons $X_b[U,V]=X[(U,0),(V,b)]$ et
$\Phi(X)=\sum_b X_b\,x^b\in M_{33}\bigl(\Z[x]/(x^3-1)\bigr)$. On a $(XY)_b=\sum_c X_cY_{b-c}$, donc $\Phi$ est un morphisme d'anneaux,
avec $\Phi(I)=I$ et $\Phi(J)=(1+x+x^2)J$. Appliquer $\Phi$ à \eqref{eq:srg}, puis évaluer en $x=1$ et en $x=\omega$
($1+\omega+\omega^2=0$), donne les équations de (i) et (ii). Les sommes de lignes de $Q$ sont les degrés, d'où $Q\one=14\one$.
On a $Q^{\mathsf T}=\sum_b A_{-b}=Q$ et $M^*=\sum_b\bar\omega^{\,b}A_{-b}=M$. Le point (iii) vient de $1-\omega^2=(1-\omega)(1+\omega)$.
Pour (iv) : $(U,0)$ est voisin de $(U,1)$ si et seulement s'il l'est de $(U,2)$, et c'est le cas exactement pour les orbites
triangulaires ; d'où $Q[U,U]=2$ ou $0$ et $M[U,U]=\omega+\omega^2=-1$ ou $0$.
\end{proof}

Le théorème~\ref{thm:libre} découle donc du lemme~\ref{lem:QM} et de l'énoncé suivant.

\begin{theorem}\label{thm:matrices}
Il n'existe pas de couple $(Q,M)$ tel que
$Q\in M_{33}(\Z)$ soit symétrique avec $Q\one=14\one$ et $Q^2+Q=12I+6J$, $M\in M_{33}(\Z[\omega])$ soit hermitienne avec $M^2+M=12I$,
et $Q\equiv M\pmod{\pi}$ coefficient par coefficient.
\end{theorem}

%====================================================================
\section{Preuve du théorème~\ref{thm:matrices}}

Soit $\rho:\Z[\omega]\to\F_3$ le morphisme de réduction ($\rho(\omega)=1$, de noyau $\pi\Z[\omega]$). Sur $\Z$, c'est la
réduction modulo~3, et $\rho(\bar z)=\rho(z)$. Pour une matrice $X$, on note $\bar X=\rho(X)$ ; alors $\rho(X^*)=\bar X^{\mathsf T}$.

\paragraph{Étape 1 (identités).}
Soit $V=\ker(Q-3I)\subset\Q^{33}$.
(a) Si $x\in V$, alors $\one^{\mathsf T}x=0$ : en effet $14\,\one^{\mathsf T}x=\one^{\mathsf T}Qx=3\,\one^{\mathsf T}x$.
(b) La matrice entière symétrique $N=11(Q+4I)-6J$ vérifie $(Q-3I)N=0$ et $Nx=77x$ pour $x\in V$.
En effet, $(Q-3I)(Q+4I)=Q^2+Q-12I=6J$ et $(Q-3I)J=11J$, donc $(Q-3I)N=66J-66J=0$ ; si $x\in V$, $(Q+4I)x=7x$ et $Jx=0$.
(c) La matrice hermitienne $N'=M+4I$ vérifie $(M-3I)N'=M^2+M-12I=0$ et $N'y=7y$ pour $y\in\ker(M-3I)$.
(d) Comme $Q$ est symétrique avec $Q\one=14\one$, on a $QJ=JQ=14J$ ; sur $\one^\perp$, $Q^2+Q=12I$, donc $14$ est valeur propre
simple et $W:=\ker\bigl((Q-3I)(Q-14I)\bigr)=V\oplus\Q\one$. La matrice entière symétrique $N_W=3(Q+4I)-J$ vérifie
$(Q-3I)(Q-14I)N_W=18(Q-14I)J=0$ et $N_Ww=21w$ pour $w\in W$ (sur $V$ par (a) ; et $N_W\one=54\one-33\one$).

\paragraph{Étape 2 (bases entières).}
Soit $R$ un anneau principal ($\Z$ ou $\Z[\omega]$) et $X\in M_n(R)$. Le module $L=\{v\in R^n : Xv=0\}$ est un facteur direct de
$R^n$ : en effet $R^n/L\simeq X(R^n)\subset R^n$ est de type fini sans torsion, donc libre. Il existe donc une matrice $B$ de taille
$n\times r$ dont les colonnes forment une base de $L$ ($r$ est la dimension du noyau sur le corps des fractions), et une matrice
$C$ de taille $r\times n$ telle que $CB=I_r$ ; alors $BCv=v$ pour tout $v\in L$. On applique ceci :
\begin{itemize}[nosep]
\item à $L=\ker_\Z(Q-3I)$ : $(B,C,r)$, $G=B^{\mathsf T}B$, $d=\det G>0$ ;
\item à $L'=\ker_{\Z[\omega]}(M-3I)$ : $(B',C',r')$, $G'=B'^*B'$, hermitienne définie positive, $d'=\det G'>0$ ;
comme $d'$ est réel et appartient à $\Z[\omega]$, $d'\in\Z$ ;
\item à $L_W=\ker_\Z\bigl((Q-3I)(Q-14I)\bigr)$, de rang $r+1$ : $(B_W,C_W)$, $G_W=B_W^{\mathsf T}B_W$.
\end{itemize}

\paragraph{Étape 3 (les identités $GK=77I$ et $G'K'=7I$).}
Posons $K=CNC^{\mathsf T}\in M_r(\Z)$. Les colonnes de $N$ sont dans $L$ (étape 1b), donc $BCN=N$ ; comme $N$ est symétrique,
$N=NC^{\mathsf T}B^{\mathsf T}$, d'où $N=BKB^{\mathsf T}$. De plus $NB=77B$, donc $B^{\mathsf T}N=77B^{\mathsf T}$ et
\[
GK=B^{\mathsf T}(BCN)C^{\mathsf T}=B^{\mathsf T}NC^{\mathsf T}=77\,(CB)^{\mathsf T}=77\,I_r .
\]
De même, avec $K'=C'N'C'^*$ et $K_W=C_WN_WC_W^{\mathsf T}$ : $N'=B'K'B'^*$, $G'K'=7I_{r'}$ et $G_WK_W=21I_{r+1}$. Conséquences :
\begin{itemize}[nosep]
\item $d\cdot\det K=77^r$, donc $d=7^a\,11^b$ ;
\item $K'$ est hermitienne, donc $\det K'\in\Z$, et $d'\cdot\det K'=7^{r'}$ ; donc $d'=7^c$, et $d'\equiv 1\pmod 3$ ;
\item $\det G_W$ divise $21^{r+1}$ ; en particulier $11\nmid\det G_W$.
\end{itemize}

\paragraph{Étape 4 (l'exposant $b$ est impair).}
Les colonnes de la matrice $[B\,|\,\one]$ sont dans $L_W$, donc $[B\,|\,\one]=B_WT$ avec $T=C_W[B\,|\,\one]\in M_{r+1}(\Z)$.
Par l'étape 1(a), $B^{\mathsf T}\one=0$ : la matrice de Gram de $[B\,|\,\one]$ vaut $\mathrm{diag}(G,33)$, et aussi
$T^{\mathsf T}G_WT$. Donc
\[
33\,d=\det(T)^2\,\det G_W .
\]
Les valuations $11$-adiques donnent $1+b=2\,v_{11}(\det T)$ : $b$ est impair. Donc $d=7^a11^b\equiv(-1)^b=-1\pmod 3$.

\paragraph{Étape 5 (réduction modulo~3 : $d\equiv d'$).}
Comme $Q\equiv M\pmod\pi$, $\bar N=2(\bar Q+I)$ et $\bar N'=\bar Q+I$ ; donc $\bar N=2\bar N'$ et $\bar N'=2\bar N$. En réduisant les
identités de l'étape~3 :
\[
\bar K=\bar C\bar N\bar C^{\mathsf T}=2\,(\bar C\bar B')\,\bar K'\,(\bar C\bar B')^{\mathsf T},\qquad
\bar K'=\bar C'\bar N'\bar C'^{\mathsf T}=2\,(\bar C'\bar B)\,\bar K\,(\bar C'\bar B)^{\mathsf T}.
\]
Les déterminants $\det K=77^r/d$ et $\det K'=7^{r'}/d'$ sont premiers à~3, donc $\bar K$ et $\bar K'$ sont inversibles. La première
identité factorise $\bar K$ par $\F_3^{\,r'}$, donc $r\le r'$ ; la seconde donne $r'\le r$. Ainsi $r=r'$, $\bar C\bar B'$ est carrée
et inversible, et en prenant les déterminants dans la première identité :
\[
2^r\,\rho(d)^{-1}=\det\bar K=2^r\,\det(\bar C\bar B')^2\,\det\bar K'=2^r\,\rho(d')^{-1},
\]
car un carré non nul de $\F_3$ vaut $1$, $\rho(77)=2$ et $\rho(7)=1$. Donc $d\equiv d'\pmod 3$.

\paragraph{Étape 6.} Les étapes 4 et 3 donnent $d\equiv 2$ et $d'\equiv 1\pmod 3$, ce qui contredit l'étape~5. \qed

\begin{remark}\label{rem:tau}
$\tr N=\tr(BKB^{\mathsf T})=\tr(KG)=77r$ et de même $\tr N'=7r'$. Avec le lemme~\ref{lem:QM}(iv), $77r=11(2\tau+132)-198$ et
$7r'=132-\tau$. L'égalité $r=r'$ obtenue à l'étape~5 donne donc $2\tau+114=132-\tau$, soit $\tau=6$ et $r=r'=18$ : c'est le
résultat d'Ishida, retrouvé par la même réduction. L'obstruction nouvelle est l'étape suivante : l'égalité des
\emph{discriminants}.
\end{remark}

%====================================================================
\section{Le cas d'un triangle fixe (proposition~\ref{prop:triangle})}

Supposons $\Fix(g)=\{t_1,t_2,t_3\}$, un triangle, et soient $(U,b)$, $U\le 32$, les sommets des orbites libres. On définit $A_b$, $Q$ et
$M$ ($32\times32$) comme au §2. On a $M^2+M=12I$ : dans le calcul du §2, les chemins passant par un $t_i$ et le terme en $J$ sont
multipliés par $1+\omega+\omega^2=0$. Et $Q\equiv M\pmod\pi$.

Les vecteurs de $\Z^{99}$ invariants par $g$ s'identifient à $\Z^{35}$ (une valeur par orbite libre, une par $t_i$). $A$ y agit
par
\[
\tilde A=\begin{pmatrix}Q&\mathbf a\\3\mathbf a^{\mathsf T}&J_3-I_3\end{pmatrix},\qquad \mathbf a[U,i]=A[(U,0),t_i],
\]
avec $\tilde A\one=14\one$. $\tilde A$ est autoadjointe pour la forme $D'=\mathrm{diag}(I_{32},\tfrac13 I_3)$. $J$ agit par
$\one w^{\mathsf T}$, avec $w=(3,\dots,3,1,1,1)$ et $w^{\mathsf T}\one=99$.

\begin{enumerate}[leftmargin=*,nosep]
\item \emph{Égalité des rangs.} $E_3=(A+4I)/7-2J/77$ est 3-entier et $E_3\equiv A+I-J\pmod 3$ s'annule sur le triangle
$\Fix(g)$. L'identité trace--rang d'Ishida \cite[cor.~4.2]{I26} donne donc $\tr\bigl(g\,|\,\ker(A-3I)\bigr)=0$. Or cette trace vaut
$r-r'$, où $r=\dim\ker(\tilde A-3I)$ et $r'=\dim\ker(M-3I)$ (un point fixe n'a pas de composante $\omega$). Donc $r=r'$.
\item \emph{Mêmes étapes.} On remplace $N$ par $N_{\mathrm{inv}}=11(\tilde A+4I)-2\,\one w^{\mathsf T}$ (qui vaut $77\,E_3$ sur les
invariants), $G$ par $\tilde G=B^{\mathsf T}D'B$ et $K$ par $K=C\,N_{\mathrm{inv}}D'^{-1}C^{\mathsf T}$, qui est entière. On vérifie
comme au §3 : les colonnes de $N_{\mathrm{inv}}$ sont dans $\ker_\Z(\tilde A-3I)$, $N_{\mathrm{inv}}B=77B$ et
$D'N_{\mathrm{inv}}D'^{-1}=N_{\mathrm{inv}}^{\mathsf T}$. On en déduit $N_{\mathrm{inv}}D'^{-1}=BKB^{\mathsf T}$ et $\tilde GK=77I$.
Modulo~3, $N_{\mathrm{inv}}D'^{-1}\equiv\mathrm{diag}\bigl(2(\bar Q+I),0\bigr)$ : les trois dernières colonnes sont multipliées par~3,
donc les trois dernières lignes aussi, par symétrie. L'étape~5, avec $r=r'$, donne alors $\det\tilde G\equiv d'\equiv1\pmod3$.
Les entrées de $\tilde G$ sont dans $\tfrac13\Z$, $\det\tilde G\cdot\det K=77^r$ et $\det K$ est premier à~3, donc $\det\tilde G$ est l'entier $77^r/\det K=7^a11^b$.
\item \emph{Le premier 11.} Avec $N_W=9(\tilde A+4I)-\one w^{\mathsf T}$ (qui vaut $63$ fois le projecteur sur
$\ker(\tilde A-3I)\oplus\Q\one$), on obtient $\tilde G_WK_W=63I$, donc $v_{11}(\det\tilde G_W)=0$. Comme
$\one^{\mathsf T}D'\one=32+1=33$, l'étape~4 donne $33\det\tilde G=\det(T)^2\det\tilde G_W$. Donc $b$ est impair et
$\det\tilde G\equiv2\pmod3$ : contradiction. \qed
\end{enumerate}

%====================================================================
\section{Remarques}

\begin{itemize}[leftmargin=*]
\item \textbf{Mécanisme.} Les deux caractères de $C_3$ voient la même forme quadratique $(\bar Q+I)$ sur $\F_3$. Seul le
caractère trivial contient le vecteur $\one$ ; sa norme $33=3\cdot11$ injecte le premier $11$, qui n'est pas un carré modulo~3.
Rien ne peut compenser : le seul autre premier en jeu est $7=r-s$, qui vaut $1$ modulo~3 (il se décompose dans $\Z[\omega]$).
\item \textbf{Ce qui est utilisé.} Pour le théorème~\ref{thm:libre} : uniquement \eqref{eq:srg} et l'action libre ; ni les valeurs
$0,1,2$ des coefficients de $Q$, ni aucun calcul. Pour la proposition~\ref{prop:triangle} : en plus, l'identité d'Ishida.
Pour le corollaire : en plus, Makhnev--Minakova.
\item \textbf{Généralisation.} La même preuve donne un critère pour tout $\srg(v,k,\lambda,\mu)$ de valeurs propres $k,\theta,\theta'$,
lorsque $3\nmid\theta-\theta'$ et que tous les premiers divisant $\theta-\theta'$ valent $1$ modulo~3 (hypothèses d'intégralité
vérifiées au cas par cas). Pour $v\le1300$, elle exclut
aussi les automorphismes d'ordre~3 sans point fixe pour $(105,52,21,30)$, $(105,52,29,22)$, $(261,52,11,10)$, $(714,248,97,80)$,
$(1221,610,279,330)$ et $(1221,610,329,280)$. Tous ces jeux sont d'existence inconnue dans les tables de Brouwer \cite{Br}.
\item \textbf{Contrôles numériques.} Le lemme « $d\equiv d'\pmod3$ » a été testé sans échec sur des graphes réels munis d'une action
libre d'ordre~3 : les graphes triangulaires $T(6)$, $T(7)$, $T(9)$, $T(10)$, $T(12)$, $T(13)$, $T(15)$, $T(18)$ et $T(21)$. Un cas réel donne $d\equiv d'\equiv2$ ($T(10)$). Il a
aussi été testé sur environ 1\,650 réseaux invariants aléatoires, avec un recalcul indépendant. Le test échoue bien lorsque le projecteur
n'est pas 3-entier, ce qui montre qu'il sait détecter une violation. La forme « $d=$ (premiers de $\theta-\theta'$) $\times$
(premiers de $\langle\one,\one\rangle$) » est observée exactement sur $T(15)$, $T(18)$ et $T(21)$.
\item \textbf{Vérification formelle.} Le théorème~\ref{thm:matrices} a été formalisé et vérifié en Lean~4
(Mathlib, version \texttt{v4.35.0-rc2}), sans \texttt{sorry} : le théorème \texttt{C3.no\_pair} ne dépend que des
axiomes standard \texttt{propext}, \texttt{Classical.choice} et \texttt{Quot.sound}. Le relecteur n'a à contrôler que
les définitions (entiers d'Eisenstein, conjugaison, réduction $\rho$) et l'énoncé. Ne sont pas encore formalisés :
le lemme~\ref{lem:QM} (passage du graphe aux matrices) et la proposition~\ref{prop:triangle}.
\item \textbf{À relire en priorité.} L'étape~2 (facteur direct sur $\Z[\omega]$), l'étape~4 (valuation 11-adique) et
l'étape~5 (réduction). Pour la proposition~\ref{prop:triangle}, la normalisation par $D'$.
\end{itemize}

{\small
\begin{thebibliography}{MM04}
\bibitem[B09]{B09} M.~Behbahani, \emph{On strongly regular graphs}, thèse, Concordia University, 2009.
\bibitem[BL11]{BL11} M.~Behbahani, C.~Lam, \emph{Strongly regular graphs with non-trivial automorphisms}, Discrete Math.\ 311 (2011).
\bibitem[Br]{Br} A.\,E.~Brouwer, \emph{Parameters of strongly regular graphs}, tables en ligne,
\url{https://aeb.win.tue.nl/graphs/srg/srgtab.html}.
\bibitem[CW23]{CW23} P.\,G.~Cesarz, A.\,J.~Woldar, \emph{On the automorphism group of a putative Conway 99-graph},
arXiv:2308.02978, 2023.
\bibitem[I26]{I26} Y.~Ishida, \emph{No involutions in the missing Moore graph}, arXiv:2606.29183v2, 2026.
\bibitem[MM04]{MM04} A.\,A.~Makhnev, I.\,M.~Minakova, Diskretnaya Matematika 16(1) (2004), 95--104.
\end{thebibliography}
}

\end{document}
